\subsection*{Question 4 (a)}

For the Toy Neural Network problem, I implemented forward finite-difference gradient descent and Nesterov random search, and compared both with exact gradient descent. In the finite-difference method, the derivative is approximated by evaluating the function after a small perturbation in each coordinate. I tested two values of the perturbation size, \(\delta=0.05\) and \(\delta=0.8\), to show how the choice of \(\delta\) affects the approximation.

\begin{figure}[H]
    \centering
    \includegraphics[width=0.6\textwidth]{q4Figures/q4b_objective_benchmarkB.png}
    \caption{Objective value versus iteration for the Toy Neural Network problem, comparing exact gradient descent, finite-difference gradient descent, and Nesterov random search.}
    \label{fig:q4b_objective_benchmarkB}
\end{figure}

Figure~\ref{fig:q4b_objective_benchmarkB} shows that exact gradient descent (blue) decreases the objective fastest and reaches a final objective value of \(0.7244\). Finite-difference gradient descent with \(\delta=0.05\) (orange) behaves very similarly, reaching \(0.7280\). This suggests that the finite-difference approximation is accurate enough when the perturbation is small. With \(\delta=0.8\) (green), the final objective is much worse, \(1.6402\), because the derivative approximation is too coarse. Nesterov random search (red) also reaches a similar final value, \(0.7254\), but takes more iterations to get there.

\begin{figure}[H]
    \centering
    \includegraphics[width=0.6\textwidth]{q4Figures/q4c_contour_benchmarkB.png}
    \caption{Contour trajectories for the Toy Neural Network problem, comparing exact gradient descent, finite-difference gradient descent, and Nesterov random search.}
    \label{fig:q4c_contour_benchmarkB}
\end{figure}

Figure~\ref{fig:q4c_contour_benchmarkB} gives the same picture from the trajectory plot. Exact gradient descent (blue) and finite-difference gradient descent with \(\delta=0.05\) (orange) move towards almost the same region. Finite-difference gradient descent with \(\delta=0.8\) (green) moves to a poorer region and then does not improve much. Nesterov random search (red) reaches the same general minimum region, but its path is less smooth because the search direction is random.

\begin{figure}[H]
    \centering
    \includegraphics[width=0.6\textwidth]{q4Figures/q4a_numerical_results.png}
    \caption{Printed numerical results for the Toy Neural Network problem.}
    \label{fig:q4a_numerical_results}
\end{figure}

Overall, the results show that finite-difference gradient descent can work well when \(\delta\) is chosen carefully, but a large \(\delta\) gives a poor approximation and harms convergence. Nesterov random search does not need exact derivatives and still finds a good solution here, although its movement is noisier than exact gradient descent.